\section{Discussion} \label{sec:discuss}

\subsection{Broader Impact}

\fbox{TODO:} some ideas:
\begin{itemize}
\item options can be used as insurance to protect against losses
\item measures can be built from the options market on the financial health and expectations of the markets, e.g.\ VIX index
\item model interpretability is important for financial implication, and particularly financial regulation. Hence controlled asymptotic make the application more meaningful.
\end{itemize}

%\subsection{Extensions}
%
%\fbox{TODO:} maybe for the conclusion?
%
%The presented approach focuses on a single stock and observations at a single time because this already represents a challenging and relevant practical problem.
%Still, we discuss how one could approach the modeling of multiple stocks, and the IVS at multiple time points.
%
%\textbf{Multi-asset.}
%It is straightforward to extend the current approach to create a model for multiple IVSs.
%Indeed, one may consider a model taking as input an input categorical variable that will the volatility surface.
%This approach would allow to transfer reuse and transfer features learned from on surface to another, which could be highly beneficial in situations where few observations are available for IVS.
%This advantage typically comes at the cost of building and training a deeper and larger NN for a joint model than for individual models.
%Note that, additional non-arbitrage constraints may have to be imposed in the specific case of implied volatility surfaces for currency pairs.
%
%\textbf{Multi-day.}
%Modeling the IVS at multiple time points could be useful to transfer information forward from the past, which may be allows to build a more resilient image of the IVS a given time point.
%This may be achieved in different ways.
%For example, the parameters $\theta$ of the previous model could be used as initial values for the model to be newly fitted.
%Alternatively, one may think of propagating forward in time a latent factor as performed in recurrent neural networks.
%
%Another application could the construction of a generative model for future IVSs which could be used for risk-management, for example.
%Such a generative model would be required to simulate entire IVSs given present and past values.
%Our approach could be used in combination with a generative module to guarantee that the fake IVSs are sensible.